\subsection{Experimental Settings}
\label{sec:setup}

\myparagraph{Platforms}
MFC in the draft sequence requires coordinated support from the HAL and image-processing components, which is currently available in Portrait mode.
We therefore conducted all experiments in Portrait mode on the Galaxy S26 Ultra and Galaxy S26, using end-to-end binaries that incorporate this support.
Table~\ref{tab:setup} summarizes these platforms and their common run conditions.
To isolate controller behavior, both evaluation binaries bypass the existing static thermal guard.
The Exynos-based Galaxy S26 disables parallel capture when either \(M\) or \(S\) is enabled, and its evaluation binary also bypasses this restriction.

\myparagraph{Control Configurations}
We compared our controller with \emph{Baseline}, which enables \(M\) and \(S\) with admission and pacing disabled.
To ensure a fair comparison, all control configurations use the same production-derived capture implementation within each experimental setting and differ only in policy.
Optional-stage skipping and pacing delays therefore reflect controller policy.

\myparagraph{Workload Conditions}
On each device, we evaluated two capture conditions: the production-default 12MP configuration under \textcolor{red}{normal conditions} and a 24MP configuration under \textcolor{red}{CPU stress}.
These are the low- and high-load endpoints of the four collected \textcolor{red}{combinations of resolution and CPU load condition}.
\textcolor{red}{For the CPU stress condition, we ran} AnTuTu Benchmark~11.1.4 concurrently with capture.
Note that 24MP denotes the requested capture resolution, because the capture pipeline falls back to 12MP at overheat level~5 and above.
For each device and capture condition, we varied the starting overheat level (0--6), resulting in 28 experimental settings.

\myparagraph{Run Protocol}
Before each run, we first reached the target starting overheat level through repeated captures or idle cooling.
We then force-stopped and relaunched the camera application to reset process-local controller state, and began the run when the device reported that level.
Thermal state then evolved without external intervention.
A reproducible ADB loop issued each shutter input as soon as the application received \texttt{captureAvailable}, saturating the exposed request rate.
Each run targeted 30 consecutive captures and terminated at the first Capture Timeout.

% Notes: docs/exhibits.md#tab_setup
\begin{table}[t]
  \centering
  \setlength{\abovecaptionskip}{2pt}%
  \setlength{\belowcaptionskip}{3pt}%
  \caption{Experimental setup.}
  \label{tab:setup}

  {\scriptsize
  \renewcommand{\arraystretch}{1.15}%
  \setlength{\tabcolsep}{3pt}%
  \begin{tabular}{>{\raggedright\arraybackslash}p{0.29\columnwidth}
                  >{\raggedright\arraybackslash}p{0.36\columnwidth}
                  >{\raggedright\arraybackslash}p{0.24\columnwidth}}
    \toprule
    \textbf{Setting} & \multicolumn{2}{c}{\textbf{Value}} \\
    \midrule
    Device & \textbf{Galaxy S26 Ultra} & \textbf{Galaxy S26} \\
    SoC & Snapdragon 8 Elite Gen 5 & Exynos~2600 \\
    RAM & \multicolumn{2}{>{\raggedright\arraybackslash}p{0.66\columnwidth}}{12\,GB} \\
    \midrule
    OS & \multicolumn{2}{>{\raggedright\arraybackslash}p{0.66\columnwidth}}{Android~17 (API level~37)} \\
    Camera software & \multicolumn{2}{>{\raggedright\arraybackslash}p{0.66\columnwidth}}{17.5.00.10, July 2026} \\
    Camera configuration & \multicolumn{2}{>{\raggedright\arraybackslash}p{0.66\columnwidth}}{Portrait mode, 1$\times$ zoom, HDR shot mode} \\
    Capture conditions & \multicolumn{2}{>{\raggedright\arraybackslash}p{0.66\columnwidth}}{12MP normal; 24MP \textcolor{red}{CPU stress}} \\
    \textcolor{red}{CPU stress} & \multicolumn{2}{>{\raggedright\arraybackslash}p{0.66\columnwidth}}{AnTuTu Benchmark~11.1.4 running during capture} \\
    Test environment & \multicolumn{2}{>{\raggedright\arraybackslash}p{0.66\columnwidth}}{Air-conditioned office, fixed scene with a static subject} \\
    Shutter input & \multicolumn{2}{>{\raggedright\arraybackslash}p{0.66\columnwidth}}{Continuous injection via an ADB loop} \\
    Run horizon & \multicolumn{2}{>{\raggedright\arraybackslash}p{0.66\columnwidth}}{Up to 30 consecutive capture requests} \\
    \bottomrule
  \end{tabular}}
\end{table}
